\section{Source-Code Interface Contract}
\label{app:source-interface-contract}

This appendix records the tensor shapes expected by the current implementation.
It is not a new mathematical layer; it is the code-facing version of the
notation introduced in the main text.%
\note{
  The source tree contains more runtime wiring than this appendix records:
  configuration binding, source selection, key-range slicing, normalization
  caches, analytics, device movement, and training wrappers.  This appendix is
  only the minimal tensor contract that the forecasting notation expects.  Some
  source modules may need to catch up as the document settles.
}

\subsection{Dataloader Batch}
\label{appsub:dataloader-batch-contract}

For one collated batch, the source dataloader exports the past window as
\[
  \text{\field{features}}
  \equiv
  \tensor{x}
  \in
  \R^{B \times C \times H_x \times D_x},
  \qquad
  \text{\field{mask}}
  \equiv
  \tensor{m}_x
  \in
  \{0,1\}^{B \times C \times H_x}.
\]
It exports the raw future window and future mask as
\[
  \text{\field{future\_features}}
  \in
  \R^{B \times C \times H_f \times D_x},
  \qquad
  \text{\field{future\_mask}}
  \equiv
  \tensor{m}_f
  \in
  \{0,1\}^{B \times C \times H_f}.
\]
The raw future width is $D_x$ because it is read from the same normalized kline
feature layout as the input window.  The MDN target-selection step then keeps
the configured future coordinates described in
Subsection~\ref{appsub:feature-registry-target-selection-contract} and forms
\[
  \tensor{f}
  \in
  \R^{B \times C \times H_f \times D_f}.
\]
Therefore $D_f$ is the selected target width used by the forecast head, not
necessarily the full raw kline width.

When key tensors are present, they follow the non-feature axes:
\[
  \text{\field{past\_keys}}
  \in
  \R^{B \times C \times H_x},
  \qquad
  \text{\field{future\_keys}}
  \in
  \R^{B \times C \times H_f}.
\]
The logical batch contract also carries the directed edge label:
\[
  \left(
    e^{b},
    \tensor{x}^{b}_{u^{b}\,v^{b}},
    {}_{x}\tensor{m}^{b}_{u^{b}\,v^{b}},
    \tensor{f}^{b}_{u^{b}\,v^{b}},
    {}_{f}\tensor{m}^{b}_{u^{b}\,v^{b}}
  \right),
  \qquad
  e^{b}=(u^{b},v^{b})\in\mathcal{E}.
\]
When this label is materialized in code, the field \field{edge\_id} is an
implementation label for the mathematical edge:
\[
  \text{\field{edge\_id}}^{b}
  \longmapsto
  e^{b}=(u^{b},v^{b})\in\mathcal{E},
  \qquad
  \text{\field{edge\_id}}\in\{1,\ldots,L\}^{B}.
\]
Some current source paths may still supply this identity through runtime binding
and source selection around the batch rather than through a materialized tensor.

In the edge-specialized setting used in this draft, the forecast module applied
to a batch is selected from the batch edge label.  Homogeneous batches satisfy
\[
  e^{1}=\cdots=e^{B}=e_{\mathrm{batch}}.
\]
The selected forecast head is then read as
\[
  {}_{\psi^{e_{\mathrm{batch}}}}g(\tensor{z}^{b}).
\]
If mixed-edge batches are introduced later, the implementation must either route
each sample by \field{edge\_id} or replace the edge-specialized heads with a
shared edge-conditioned head.

\subsection{NodeLift Graph-Collated Interface}
\label{appsub:nodelift-graph-collated-interface}

NodeLift is a graph-level transformation, not an edge-local batch operation.  It
requires a usable observed input edge field over several edges at the same
anchor time, channel, and position.  Therefore an implementation that starts
from edge-labeled samples must provide a graph-collation step.

For one graph-collated observed-input anchor batch, a convenient code-facing
shape is
\[
  \text{\field{edge\_features}}
  \in
  \R^{B_g\times L_{\mathrm{use}}\times C\times H_x\times D_x},
\]
where $B_g\in\N$ is the number of graph anchors in the graph-collated batch.
The corresponding edge mask is
\[
  \text{\field{edge\_mask}}
  \in
  \{0,1\}^{B_g\times L_{\mathrm{use}}\times C\times H_x}.
\]
The coordinatewise edge mask is
\[
  \text{\field{edge\_coord\_mask}}
  \in
  \{0,1\}^{B_g\times L_{\mathrm{use}}\times C\times H_x\times D_y}.
\]
The graph metadata provides
\[
  \text{\field{edge\_index}}
  \longmapsto
  e=(u,v),
  \qquad
  \mat{A},
  \qquad
  \mat{B},
  \qquad
  \mat{Q},
  \qquad
  \mat{U},
  \qquad
  \mat{G}_{\star}.
\]

NodeLift returns
\[
  \text{\field{node\_features}}
  \in
  \R^{B_g\times C\times H_x\times N\times D_y},
\]
and the primary coordinatewise node mask
\[
  \text{\field{node\_mask}}
  \in
  \{0,1\}^{B_g\times C\times H_x\times N\times D_y}.
\]
A coarse node mask may also be exported when needed:
\[
  \text{\field{node\_mask\_any}}
  \in
  \{0,1\}^{B_g\times C\times H_x\times N},
\]
or
\[
  \text{\field{node\_mask\_all}}
  \in
  \{0,1\}^{B_g\times C\times H_x\times N}.
\]
The price residual sidecar is
\[
  \text{\field{price\_residual}}
  \in
  \R^{B_g\times C\times H_x\times L_{\mathrm{use}}\times
  |\mathcal{J}_{\mathrm{price}}|}.
\]
Its residual mask is
\[
  \text{\field{price\_residual\_mask}}
  \in
  \{0,1\}^{B_g\times C\times H_x\times L_{\mathrm{use}}\times
  |\mathcal{J}_{\mathrm{price}}|}.
\]
When activity diagnostics are retained, NodeLift may also return
\[
  \text{\field{activity\_support}},
  \qquad
  \text{\field{activity\_coverage}},
  \qquad
  \text{\field{activity\_total}},
\]
corresponding to $\vect{\kappa}^{c\,i\,d}$, $\rho_u^{c\,i\,d}$, and
$\vect{\nu}_{\Sigma}^{c\,i\,d}$.  Their common code-facing shape is
\[
  B_g\times C\times H_x\times N\times |\mathcal{J}_{\mathrm{act}}|.
\]

The important requirement is that the observed edge input field is aligned over
a common graph anchor before NodeLift is applied.

\subsection{Feature Registry and Target Selection}
\label{appsub:feature-registry-target-selection-contract}

The normalized kline source registry has $D_x=9$ model coordinates:
\begin{center}
\begin{tabular}{c l l}
\toprule
Coordinate & Field & Meaning \\
\midrule
1 & \field{open} & Log open relative to previous close. \\
2 & \field{high} & Log high relative to previous close. \\
3 & \field{low} & Log low relative to previous close. \\
4 & \field{close} & Log close relative to previous close. \\
5 & \field{volume} & $\log(1+\cdot)$ applied to volume. \\
6 & \field{quote\_volume} & $\log(1+\cdot)$ applied to quote volume. \\
7 & \field{trades} & $\log(1+\cdot)$ applied to number of trades. \\
8 & \field{taker\_buy\_base} & $\log(1+\cdot)$ applied to taker buy base volume. \\
9 & \field{taker\_buy\_quote} & $\log(1+\cdot)$ applied to taker buy quote volume. \\
\bottomrule
\end{tabular}
\end{center}

The target feature registry is an ordered selection from this source registry:
\[
  \mathcal{J}_f
  =
  (d_1,\ldots,d_{D_f}),
  \qquad
  d_a\in\{1,\ldots,D_x\},
  \qquad
  a=1,\ldots,D_f.
\]
Thus $D_f$ is the number of selected future coordinates, while
$\mathcal{J}_f$ is the selected coordinate list.  The target tensor is formed by
\[
  f^{b\,c\,i\,a}
  =
  \text{\field{future\_features}}^{b\,c\,i\,d_a},
  \qquad
  a=1,\ldots,D_f.
\]

The scalar edge-return readout uses the close-return source coordinate
\[
  d_{(r)}=4.
\]
When scalar edge-return recovery is enabled, $\mathcal{J}_f$ must contain this
coordinate, and the label $(r)$ denotes the selected target slot corresponding
to $d_{(r)}$.

\subsection{Representation Output}
\label{appsub:representation-output-contract}

The node-first representation encoder consumes the NodeLift output
\[
  (\text{\field{node\_features}},\text{\field{node\_mask}})
  \in
  \R^{B_g \times C \times H_x \times N \times D_y}
  \times
  \{0,1\}^{B_g \times C \times H_x \times N \times D_y}
\]
or an equivalent flattened batch of node windows.  It returns node encodings
\[
  \text{\field{node\_encoding}}
  \equiv
  \tensor{z}_{\mathcal{V}}
  \in
  \R^{B_g\times N\times D_e}.
\]
For an edge-specialized MDN, code then builds an edge-conditioned encoding
\[
  \text{\field{encoding}}
  \equiv
  \tensor{z}_{\mathcal{E}}
  \in
  \R^{B_g\times L_{\mathrm{use}}\times D_e}
\]
from endpoint node encodings and optional residual sidecars.  A pure MDN call on
a single edge batch may still receive a flattened
$\text{\field{encoding}}\in\R^{B\times D_e}$.  If an encoder path produces a
sequence encoding, it must be reduced before entering the pure MDN model.

\subsection{MDN Output}
\label{appsub:mdn-output-contract}

The current MDN implementation takes an edge-conditioned encoding
\[
  \text{\field{encoding}}
  \in
  \R^{B \times D_e}
\]
and returns the structure \field{MdnOut}.  Its fields are
\[
  \text{\field{log\_pi}}
  \in
  \R^{B \times C \times H_f \times K},
\]
\[
  \text{\field{mu}}
  \in
  \R^{B \times C \times H_f \times K \times D_f},
  \qquad
  \text{\field{sigma}}
  \in
  \Rpos^{B \times C \times H_f \times K \times D_f}.
\]
The \field{log\_pi} tensor contains log-normalized mixture weights.  Thus
$\exp(\text{\field{log\_pi}})$ corresponds to the mixture coefficient tensor
$\tensor{\alpha}$.  The \field{sigma} tensor is already positive and represents
the diagonal component scale vector used to form component covariances.

\subsection{Loss Interface}
\label{appsub:loss-interface-contract}

The MDN negative-log-likelihood reads
\[
  \text{\field{MdnOut}},
  \qquad
  \tensor{f}
  \in
  \R^{B \times C \times H_f \times D_f},
  \qquad
  \tensor{m}_f
  \in
  \{0,1\}^{B \times C \times H_f}.
\]
The unreduced NLL map has one scalar per batch, channel, and horizon:
\[
  \ell_{\mathrm{map}}
  \in
  \R^{B \times C \times H_f}.
\]
Only positions with $m_f^{b\,c\,i}=1$ contribute to the reduced loss.  Optional
channel, horizon, and feature weights may be applied before reduction; the
reduced training loss is a scalar in $\R$.
